%% toc: P15 161 161.3,161.12
\chapter{Observation Fields Smaller than the Sample Space}
\label{sec:p15-w02c-p15-017-observation-fields}

\Cref{def:p15-tail-sigma-algebra} places the tail field \(\mathcal T\)
inside the field \(\mathcal H_\infty=\sigma(X_1,X_2,\ldots)\) generated by
the whole observed sequence, and adds that \(\mathcal H_\infty\) need not
be the ambient field \(\mathcal F\) of the probability space. This chapter
proves that boundary. It gives a construction that attaches to any
sequence an event lying outside \(\mathcal H_\infty\), and shows that the
gap is not removable by ignoring null sets. It then records the one
common setting in which the gap is absent, the canonical coordinate
product.

\emph{Inputs.} A probability space \((\Omega',\mathcal F',P')\) carrying
random elements \(X'_n\) with values in measurable spaces
\((S_n,\mathcal S_n)\), \(n\geq1\); or a nonempty measurable space
\((S,\mathcal S)\) for the canonical case.
\emph{Conclusions.} An enlarged probability space on which the same
sequence generates a proper sub-sigma-algebra of the ambient field, an
event at distance \(1/2\) in probability from every event of
\(\mathcal H_\infty\), and the equality \(\mathcal H_\infty=\mathcal F\)
for coordinate maps on a canonical product.

\section{An event the sequence does not observe}
\label{sec:p15-w02c-p15-017-unobserved-event}

\begin{hbdefinition}{def:p15-w02c-p15-017-coin-enlargement}{An unrecorded fair coin}
Let \((\Omega',\mathcal F',P')\) be a probability space and let
\(X'_n:\Omega'\to S_n\) be measurable for each \(n\geq1\). Put
\[
 \Omega=\Omega'\times\{0,1\},\qquad
 \mathcal F=\mathcal F'\otimes2^{\{0,1\}},\qquad
 P=P'\otimes\bigl(\tfrac12\delta_0+\tfrac12\delta_1\bigr),
\]
where \(2^{\{0,1\}}\) is the power set and \(\delta_y\) is the unit mass
at \(y\). Define \(X_n(\omega',y)=X'_n(\omega')\) and
\(Y(\omega',y)=y\). The fields \(\mathcal H_n\), \(\mathcal H_\infty\),
\(\mathcal F'_n\) and \(\mathcal T\) of
\cref{def:p15-tail-sigma-algebra} are formed from \(X_1,X_2,\ldots\) on
\((\Omega,\mathcal F,P)\). Write
\[
 \mathcal G=\{A'\times\{0,1\}:A'\in\mathcal F'\},\qquad
 B=\{Y=1\}=\Omega'\times\{1\}.
\]
\end{hbdefinition}

Here \(P\) is the product of two probability measures on a measurable
rectangle field. On a rectangle \(A'\times D\) its value is
\(P'(A')\,|D|/2\), and only such values are used below.

\begin{hbtheorem}{thm:p15-w02c-p15-017-observation-field-proper}{proved=annex:p15-w02c-p15-017-observation-field-proper}{The observed field can be a proper part of the ambient field}
In \cref{def:p15-w02c-p15-017-coin-enlargement}:
\begin{enumerate}
\item \(\mathcal G\) is a sigma-algebra contained in \(\mathcal F\), and
\(\mathcal T\subseteq\mathcal H_\infty\subseteq\mathcal G\).
\item \(B\in\mathcal F\) but \(B\notin\mathcal G\). Hence
\(\mathcal H_\infty\neq\mathcal F\), whatever the sequence
\((X'_n)\) is.
\item For every \(A\in\mathcal H_\infty\),
\(P(A\triangle B)=1/2\) and \(P(A\cap B)=P(A)P(B)\).
In particular no event of \(\mathcal H_\infty\) agrees with \(B\)
up to a null set, so \(\mathcal H_\infty\) differs from \(\mathcal F\)
even after events are identified when their symmetric difference is null.
\item For every \(n\) and every \(C\in\mathcal H_n\),
\(P(B\triangle C)=1/2\). The approximation of
\cref{thm:p15-finite-coordinate-approximation} therefore fails for
\(B\); its hypothesis \(A\in\mathcal H_\infty\) cannot be weakened to
\(A\in\mathcal F\).
\end{enumerate}
\end{hbtheorem}

\begin{hbexample}{def:p15-w02c-p15-017-coin-example}{A recorded sequence and one unrecorded toss}
Take \(\Omega'=\{0,1\}^{\mathbb N_+}\) with its coordinate product field
and the product of fair coin laws, and let \(X'_n\) be the \(n\)th
coordinate. On the enlarged space the sequence \(X_1,X_2,\ldots\)
records every toss but one, and \(Y\) is that remaining toss. The event
\(\{Y=1\}\) has probability \(1/2\). Every event determined by the
recorded tosses, for instance \(\{X_n=1\text{ i.o.}\}\) or
\(\{X_1=1\}\), differs from \(\{Y=1\}\) on a set of probability exactly
\(1/2\). The independent sequence obeys
\cref{thm:p15-kolmogorov-tail-zero-one}, so \(\{X_n=1\text{ i.o.}\}\)
has probability zero or one; it has probability one by
\cref{thm:p15-borel-cantelli-second}, since the events
\(\{X_n=1\}\) are mutually independent with probabilities summing to
infinity. By part~3 of the theorem, \(\{Y=1\}\) is independent of all
of \(\mathcal H_\infty\), and its probability is \(1/2\).

The smallest case takes \(\Omega'\) to be a single point, with every
\(X'_n\) constant. Then \(\Omega\) has two points,
\(\mathcal H_\infty=\mathcal T=\{\emptyset,\Omega\}\), and
\(\mathcal F\) has four events.
\end{hbexample}

\begin{hbwarning}{def:p15-w02c-p15-017-independence-not-zero-one}{Independence from the observed field does not force zero or one}
The proof of \cref{thm:p15-kolmogorov-tail-zero-one} uses two facts
about a tail event \(A\): it is independent of \(\mathcal H_\infty\),
and it belongs to \(\mathcal H_\infty\). Then \(A\) is independent of
itself, so \(P(A)=P(A)^2\). Durrett's proof of his Theorem~2.5.3 is
organized in the same way: his step (b) gives independence from
\(\sigma(X_1,X_2,\ldots)\), and the inclusion of \(\mathcal T\) in that
field then makes a tail event independent of itself
\cite[pp.~81--82]{DurrettPTEAuthor2019}.
The event \(B\) above has the first property
but not the second, and \(P(B)=1/2\). The inclusion
\(\mathcal T\subseteq\mathcal H_\infty\) in
\cref{def:p15-tail-sigma-algebra} is therefore a working premise of the
zero-one conclusion, not a formality.
\end{hbwarning}

\section{The canonical case}
\label{sec:p15-w02c-p15-017-canonical}

\begin{hbtheorem}{thm:p15-w02c-p15-017-canonical-generated}{proved=annex:p15-w02c-p15-017-canonical-generated}{Coordinate maps generate the canonical product field}
Let \((S,\mathcal S)\) be a nonempty measurable space,
\(\Omega=S^{\mathbb N_+}\), and let \(\mathcal F\) be the coordinate
product sigma-algebra, the sigma-algebra generated by the sets
\(\{\omega:\omega_i\in C\}\) with \(i\geq1\) and \(C\in\mathcal S\).
For the coordinate maps \(X_i(\omega)=\omega_i\),
\(\mathcal H_\infty=\mathcal F\).
The same holds with \(S\) replaced by a different measurable space at
each coordinate.
\end{hbtheorem}

The two theorems separate a property of a model from a property of a
sequence. On the canonical product, every event is an event about the
observed coordinates, which is the setting of
\cref{def:p15-permutable-events} and of the finite-change criterion in
\cref{thm:p15-tail-event-structure}. On a general probability space the
same sequence may leave events unobserved. The zero-one laws and the
finite-coordinate approximation speak only about events inside
\(\mathcal H_\infty\).

\emph{Sources.} The fields involved are those of
\cref{def:p15-tail-sigma-algebra}; Durrett introduces the future fields
\(\mathcal F'_n\) and the tail field \(\mathcal T\) as their intersection
at \cite[p.~81]{DurrettPTEAuthor2019}. The enlargement construction and
the proofs of both theorems are elementary and are given in full in the
annex; they are not taken from a source.
